\documentclass[tikz,border=4pt]{standalone}
\usepackage{ctex}
\usetikzlibrary{positioning}
\begin{document}
\begin{tikzpicture}[
    box/.style={rectangle, draw, rounded corners=4pt, minimum width=5.5cm, minimum height=1cm, font=\small, align=center, fill=blue!5},
    title/.style={font=\bfseries\large},
]
\node[title] (t) at (0,0) {反常积分 $p$-判别表};

\node[box, below left=0.6cm and 0.4cm of t, fill=green!10] (a) {**$\int_1^{+\infty}\frac{1}{x^p}\,dx$**\\$p>1$ → **收敛**\\$p\leq 1$ → **发散**};

\node[box, below right=0.6cm and 0.4cm of t, fill=red!10] (b) {**$\int_0^1\frac{1}{x^p}\,dx$**\\$p<1$ → **收敛**\\$p\geq 1$ → **发散**};

\node[font=\bfseries, below=4cm of t] {⚠ 临界值 $p=1$ 两边相反！};

\node[font=\small, below=4.7cm of t, text width=10cm, align=center]
  {**记忆**：无穷尾部要"衰减得足够快"（$p$ 大）；\\奇点附近要"奇异性足够弱"（$p$ 小）。\\$p=1$ 是 $\ln$ 型，恰好处于临界（发散）。};
\end{tikzpicture}
\end{document}
